\section{Introduction}\label{sec-intro}

\paragraph{Contexte.}
La République du Bénin a engagé la mise en place d'un système national
de vidéoprotection couvrant cinq villes et les postes
frontaliers, décidée en Conseils des Ministres du 11 mars et du 3 juin
2026~\cite{sgg2026conseil}. Au-delà de la sécurité publique, un tel
réseau de caméras ouvre la voie à la vidéo-verbalisation : un pipeline
de vision par ordinateur qui enchaîne détection des véhicules, suivi
multi-objets, détection de plaque d'immatriculation, reconnaissance
optique de caractères, application de règles d'infraction et restitution
sur tableau de bord (\cref{fig-pipeline}).

\begin{figure}[b]
  \centering
  \begin{tikzpicture}[
      maillon/.style={draw=bleunuit!70, rounded corners=2pt, align=center,
                      font=\footnotesize, minimum height=10mm,
                      inner sep=3pt, text width=20mm},
      vise/.style={draw=terrecuite, thick, fill=terrecuite!10},
      fleche/.style={->, >={Stealth[length=2mm]}, semithick,
                     draw=bleunuit!60},
      node distance=3.5mm]
    \node[maillon, vise]        (det)   {Détection\\des véhicules};
    \node[maillon, right=of det]  (suivi) {Suivi\\multi-objets};
    \node[maillon, right=of suivi] (plaq) {Détection\\de plaques};
    \node[maillon, right=of plaq]  (ocr)  {Lecture\\optique (OCR)};
    \node[maillon, right=of ocr]   (regles) {Règles\\d'infraction};
    \node[maillon, right=of regles] (bord) {Tableau\\de bord};
    \draw[fleche] (det)    -- (suivi);
    \draw[fleche] (suivi)  -- (plaq);
    \draw[fleche] (plaq)   -- (ocr);
    \draw[fleche] (ocr)    -- (regles);
    \draw[fleche] (regles) -- (bord);
    \node[below=1.2mm of det, font=\scriptsize\itshape,
          text=terrecuite] {ce travail};
  \end{tikzpicture}
  \caption{La chaîne de vidéo-verbalisation, du flux caméra au tableau
  de bord. Ce travail porte sur le premier maillon : un véhicule non
  détecté est perdu pour toute la suite de la chaîne.}
  \label{fig-pipeline}
\end{figure}

\paragraph{Problème.}
Tout ce pipeline repose sur son premier maillon : si un véhicule n'est
pas détecté, rien en aval ne peut le rattraper. Or les détecteurs
génériques, pré-entraînés sur des jeux comme COCO~\cite{lin2014coco},
échouent précisément dans les configurations les plus fréquentes du
trafic urbain ouest-africain : véhicules éloignés de la caméra donc de
petite taille à l'image, et véhicules partiellement masqués dans les
embouteillages. Les auteurs de BMD-45~\cite{sharma2026bmd45} quantifient
cet écart de domaine : des détecteurs entraînés sur le benchmark
UA-DETRAC~\cite{wen2020uadetrac} plafonnent à $33{,}6\,\%$ de
mAP@0,5:0,95 sur du trafic urbain dense de pays en développement, contre
$83{,}8\,\%$ pour des modèles entraînés en domaine, soit un facteur
$2{,}5$.

\paragraph{Périmètre et question de recherche.}
Ce travail se restreint volontairement au premier maillon, la détection
de véhicules, et à l'intérieur de celui-ci à un verrou précis : les
véhicules éloignés, de petite taille ou masqués. Le suivi multi-objets,
la lecture de plaques, les règles d'infraction et l'estimation de
vitesse sont hors périmètre. La question de recherche est la suivante :
\emph{un fine-tuning de YOLOv8 sur du trafic urbain dense améliore-t-il
la détection des véhicules, et spécifiquement celle des véhicules de
petite taille ?} Aucun jeu de données public de vision par ordinateur
n'existant pour le Bénin, l'étude s'appuie sur BMD-45, un jeu de trafic
urbain indien structurellement comparable (caméras CCTV fixes, forte
densité, dominance des deux-roues), sans prétendre pour autant à une
adaptation au trafic béninois lui-même (\cref{sec-discussion}).

\paragraph{Contributions.}
Nos contributions sont au nombre de trois :
\begin{enumerate}
  \item un protocole de comparaison équitable entre un modèle
        pré-entraîné (80 classes COCO) et un modèle fine-tuné
        (4 classes), fondé sur le remappage des annotations vers un
        espace de classes aligné sur COCO et leur réindexation à la
        volée, de sorte que les deux modèles sont évalués sur
        exactement les mêmes images et les mêmes boîtes de référence
        (\cref{sec-protocole}) ;
  \item le fine-tuning reproductible d'un YOLOv8n sur un sous-ensemble
        de 3\,000 images de BMD-45 extrait avec graines fixées, avec des
        augmentations choisies pour la problématique des petits objets
        (\cref{sec-entrainement}) ;
  \item une évaluation ventilée par taille d'objet selon la convention
        COCO, avec normalisation des aires à la résolution d'entrée du
        réseau, qui montre que le gain de rappel est d'autant plus fort
        que les objets sont petits : de $0{,}122$ à $0{,}694$
        ($+469{,}4\,\%$) pour les objets de moins de $32 \times 32$
        pixels (\cref{sec-resultats}).
\end{enumerate}

La suite de l'article est organisée ainsi : la \cref{sec-related} situe
le travail, la \cref{sec-methode} détaille données, entraînement et
protocole d'évaluation, la \cref{sec-resultats} rapporte les résultats,
la \cref{sec-discussion} les interprète et en énonce les limites, et la
\cref{sec-conclusion} conclut.
